% ==== AV ====
\yoschapterstar{Appendix to Chapter V. Weak Topologies and Duality in Locally Convex Linear Topological Spaces}
\label{chap:AV}

The present book is so designed that the reader may skip this appendix in the first reading and proceed directly to the following chapters. 


\yossection{1}{Polar Sets}
\label{sec:AV.1}

\begin{definitionu} Let $X$ be a locally convex linear topological space. For any set $M \subseteq X$, we define its (right) polar set $M^0$ by 
\begin{equation}\label{eq:AV.1.1}
M ^{0} = \left\{x ^{\prime} \in X ^{\prime}; \sup\limits _{x \in M} | \langle x, x ^{\prime} \rangle | \leqq 1 \right\}.
\end{equation}
\end{definitionu}

Similarly, for any set $M' \subseteq X'$, we define its (left) polar set $^0M'$ by 
\begin{equation}\label{eq:AV.1.2}
{ } ^{0} M ^{ \prime } = \{ x \in X ; \operatorname* { s u p } _ { x ^{ \prime } \in M ^{ \prime } } | \langle x , x ^{ \prime } \rangle | \leqq 1 \} = X \cap ( M ^{ \prime } ) ^{0} ,
\end{equation}
where we consider $X$ to be embedded in its bidual $(X_s')'$. 

A fundamental system of neighbourhoods of $0$ in the weak topology of $X$ is given by the system of sets of the form $^{0}M'$ where $M'$ ranges over arbitrary finite sets of $X'$. A fundamental system of neighbourhoods of $0$ in the weak* topology of $X'$ is given by the system of sets of the form $M^{0}$ where $M$ ranges over arbitrary finite sets of $X$. A fundamental system of neighbourhoods of $0$ in the strong topology of $X'$ is given by the system of sets of the form $M^{0}$ where $M$ ranges over arbitrary bounded sets of $X$. 

\begin{propositionu} $M^{0}$ is a convex, balanced set closed in the weak* topology of $X'$. 
\end{propositionu}

\begin{proof}
 For any fixed $x \in X$, the linear functional $f(x') = \langle x, x' \rangle$ is continuous in the weak* topology of $X'$. Thus $M^0 = \bigcap\limits_{m \in M} \{m\}^0$ is closed in the weak* topology of $X'$. The balanced convexity of $M^0$ is clear. 
\end{proof}

\yossub{An Application of Tychonov's Theorem}

\begin{theorem}\label{thm:AV.1.Theorem.1} Let $X$ be a locally convex linear topological space, and $A$ a convex, balanced neighbourhood of $0$ of $X$. Then $A^{0}$ is compact in the weak* topology of $X'$. 
\end{theorem}

\begin{proof}
 Let $p(x)$ be the Minkowski functional of $A$. Consider, for each $x \in X$, a sphere $S_x = \{z \in C; |z| \leqq p(x)\}$ and the topological product $S = \prod\limits_{x \in X} S_x$. $S$ is compact by Tychonov's theorem. Any element $x' \in X'$ is determined by the set of values $x'(x) = \langle x, x' \rangle$, $x \in X$. Since $x \in (p(x) + \varepsilon)A$ for any $\varepsilon > 0$, we see that $x' \in X'$ implies $\langle x, x' \rangle = \langle (p(x) + \varepsilon)a, x' \rangle$ with a certain $a \in A$. Thus $x \in A^0$ implies that $|x'(x)| \leqq p(x) + \varepsilon$, that is, $x'(x) \in S_x$. Hence we may consider $A^0$ as a subset of $S$. Moreover, it is easy to verify that the topology induced on $A^0$ by the weak* topology of $X'$ is the same as the topology induced on $A^0$ in the Cartesian product topology of $S = \prod\limits_{x \in X} S_x$. 

Hence it is sufficient to prove that $A^0$ is a closed subset of $S$. Suppose $y = \prod\limits_{x \in X} y(x)$ is an element of the weak* closure of $A^0$ in $S$. Consider any $\varepsilon > 0$ and any $x_1, x_2 \in X$. The set of all $u = \prod\limits_{x \in X} u(x) \in S$ such that 
$|u(x_{1})-y(x_{1})|<\varepsilon, |u(x_{2})-y(x_{2})|<\varepsilon$ and $|u(x_{1}+x_{2})-y(x_{1}+x_{2})|<\varepsilon$ is a neighbourhood of $y$ in $S$. This neighbourhood contains some point $x'\in A^{0}$ and, since $x'$ is a continuous linear functional on $X$, we have 
\[
\begin{aligned}
\left| y (x _{1} + x _{2}) - y (x _{1}) - y (x _{2}) \right|& \leqq \left| y (x _{1} + x _{2}) - \langle x _{1} + x _{2}, x ^{\prime} \rangle \right|
\\
\quad + \left| \langle x _{1}, x ^{\prime} \rangle - y (x _{1}) \right| + \left| \langle x _{2}, x ^{\prime} \rangle - y (x _{2}) \right|& < 3 \varepsilon .
\end{aligned}
\]

This proves that $y(x_1 + x_2) = y(x_1) + y(x_2)$. Similarly, we prove that $y(\beta x) = \beta y(x)$, and so $y$ defines a linear functional on $X$. By the fact that $y = \prod\limits_{x \in X} y(x) \in S$, we know that $|y(x)| \leqq p(x)$. Since $p(x)$ is continuous, $y(x)$ is a continuous linear functional, i.e., $y \in X'$. On the other hand, since $y$ is a weak* accumulation point of $A^0$, there exists, for any $\varepsilon > 0$ and $a \in A$, an $x' \in A^0$ such that $|y(a) - \langle a, x' \rangle| \leqq \varepsilon$. Hence $|y(a)| \leqq |\langle a, x' \rangle| + \varepsilon \leqq 1 + \varepsilon$, and so $|y(a)| \leqq 1$, that is, $y \in A^0$. 
\end{proof}

\begin{corollaryu} The unit sphere $S^{*} = \{x' \in X'; \| x' \| \leqq 1\}$ of the dual space $X_{s}'$ of a normed linear space $X$ is compact in the weak* topology of $X'$. 
\end{corollaryu}

\yossub{An Application of Mazur's Theorem}

\begin{theorem}\label{thm:AV.1.Theorem.2} Let $M$ be a convex, balanced closed set of a locally convex linear topological space $X$. Then $M = {}^0 (M^0)$. 
\end{theorem}

\begin{proof}
 It is clear that $M \subseteq {}^0 (M^0)$. If there exists an $x_0 \in {}^0 (M^0) - M$, then, by Mazur's theorem $3$ in Chapter~\ref{chap:IV}, \ref{sec:IV.6}, there exists an $x_0' \in X'$ such that $\langle x_0, x_0' \rangle > 1$ and $|\langle x, x_0' \rangle| \leqq 1$ for all $x \in M$. The last inequality shows that $x_0' \in M^0$ and so $x_0$ cannot belong to ${}^0 (M^0)$. 
\end{proof}


\yossection{2}{Barrel Spaces}
\label{sec:AV.2}

\begin{definitionu} In a locally convex linear topological space $X$, any convex, balanced and absorbing closed set is called a barrel (tonneau in Bourbaki's terminology). $X$ is called a barrel space if each of its barrels is a neighbourhood of $0$. 
\end{definitionu}

\begin{theorem}\label{thm:AV.2.Theorem.1} $A$ locally convex linear topological space $X$ is a barrel space if $X$ is not of the first category. 
\end{theorem}

\begin{proof}
 Let $T$ be a barrel in $X$. Since $T$ is absorbing, $X$ is the union of closed sets $nT = \{nt; t \in T\}$, where $n$ runs over positive integers. Since $X$ is not of the first category, at least one of the $(nT)'s$ contains an interior point. Hence $T$ itself contains an interior point $x_{0}$. If $x_{0} = 0$, $T$ is a neighbourhood of $0$. If $x_{0} \neq 0$, then $-x_{0} \in T$ by the fact that $T$ is balanced. Thus $-x_{0}$ is an interior point of $T$ with $x_{0}$. This proves that the convex set $T$ contains $0 = (x_{0} - x_{0})/2$ as an interior point. 
\end{proof}

\begin{corollary}\label{cor:AV.2.Corollary.1} All locally convex $F$-spaces and, in particular, all $B$-spaces and $\mathfrak{E}(\mathbb{R}^{n})$ are barrel spaces. 
\end{corollary}

\begin{corollary}\label{cor:AV.2.Corollary.2} The metric linear space $\mathfrak{D}_K(\mathbb{R}^{n})$ is a barrel space. 
\end{corollary}

\begin{proof}
 Let $\{\varphi_{k}\}$ be a Cauchy sequence with respect to the distance 
\[
\dist (\varphi , \psi) = \sum\limits_{m=0} ^{\infty} 2 ^{-m} \frac {p _{m} (\varphi - \psi)}{1 + p _{m} (\varphi - \psi)}, \text{ where } p _{m} (\varphi) = \sup\limits _{| j | \leqq m, x \in K} | D ^{j} \varphi (x) |. \]

For any differential operator $D^{j}$, the sequence $\{D^{j}\varphi_{k}(x)\}$ is equi-continuous and equi-bounded, that is, 
\[
\lim\limits _ {| x ^{1} - x ^{2} | \downarrow 0} \sup\limits _{k \geqq 1} \left| D ^{j} \varphi_{k} (x ^{1}) - D ^{j} \varphi_{k} (x ^{2}) \right| = 0 \text{ and } \sup\limits _{x \in K, k \geqq 1} \left| D ^{j} \varphi_{k} (x) \right| < \infty .
\]

This we see from the fact that, for any coordinate $x_s$, $\sup\limits_{x \in K, k \geqq 1} \left| \frac{\partial}{\partial x_s} D^j \varphi_k(x) \right| < \infty$. Hence, by the Ascoli-Arzela theorem, there exists a subsequence of $\{D^j \varphi_{k'}(x)\}$ which converges uniformly on $K$. By the diagonal method, we may choose a subsequence $\{\varphi_{k''}(x)\}$ of $\{\varphi_k(x)\}$ such that, for any differential operator $D^j$, the sequence $\{D^j \varphi_{k''}(x)\}$ converges uniformly on $K$. Thus 
\[
\lim\limits _ {k ^{\prime \prime} \rightarrow \infty} D ^{j} \varphi_ {k ^{\prime \prime}} (x) = D ^{j} \varphi (x) \text{ where } \varphi (x) = \lim\limits _ {k ^{\prime \prime} \rightarrow \infty} \varphi_ {k ^{\prime \prime}} (x),
\]
and these limit relations hold uniformly on $K$. Hence the metric space $\mathfrak{D}_K(\mathbb{R}^{n})$ is complete and so it is not of the first category. 
\end{proof}

\begin{remarku} \begin{enumerate}[(i)]
\item The above proof shows that a bounded set of $\mathfrak{D}(\mathbb{R}^{n})$ is relatively compact in the topology of $\mathfrak{D}(\mathbb{R}^{n})$. For, a bounded set $B$ of $\mathfrak{D}(\mathbb{R}^{n})$ is contained in some $\mathfrak{D}_K(\mathbb{R}^{n})$ where $K$ is a compact set of $R_n^n$, and, moreover, the boundedness condition of $B$ implies the equi-boundedness and equi-continuity of $\{D^j\varphi; \varphi \in B\}$ for every $D^j$.
\item Similarly, we see that any bounded set of $\mathfrak{E}(\mathbb{R}^{n})$ is a relatively compact set of $\mathfrak{E}(\mathbb{R}^{n})$.
\end{enumerate}
\end{remarku}

\begin{corollary}\label{cor:AV.2.Corollary.3} $\mathfrak{D}(\mathbb{R}^{n})$ is a barrel space. 
\end{corollary}

\begin{proof}
 $\mathfrak{D}(\mathbb{R}^{n})$ being an inductive limit of $\{\mathfrak{D}_K(\mathbb{R}^{n})\}$ when $K$ ranges over compact subsets of $\mathbb{R}^{n}$, the present corollary is a consequence of the following 
\end{proof}

\begin{propositionu} Let a locally convex linear topological space $X$ be an inductive limit of its barrel subspaces $X_{\alpha}, \alpha \in A$. Then $X$ itself is a barrel space. 
\end{propositionu}

\begin{proof}
 Let $V$ be a barrel of $X$. By the continuity of the identical mapping $T_{\alpha}: x \to x$ of $X_{\alpha}$ into $X$, the inverse image $T_{\alpha}^{-1}(V) = V \cap X_{\alpha}$ is closed with $V$. Thus $V \cap X_{\alpha}$ is a barrel of $X_{\alpha}$. $X_{\alpha}$ being a barrel space, $V \cap X_{\alpha}$ is a neighbourhood of $0$ of $X_{\alpha}$. $X$ being an inductive limit of $X_{\alpha}$'s, $V$ must be a neighbourhood of $0$ of $X$. 
\end{proof}

\begin{theorem}\label{thm:AV.2.Theorem.2} Let $X$ be a barrel space. Then the mapping $x \to Jx$ of $X$ into $(X_{s}^{\prime})_{s}^{\prime}$, defined in Chapter~\ref{chap:IV}, \ref{sec:IV.8}, is a topological mapping of $X$ onto JX, where the topology of JX is provided with the relative topology of JX as a subset of $(X_{s}^{\prime})_{s}^{\prime}$. 
\end{theorem}

\begin{proof}
 Let $B'$ be a bounded set of $X_s'$. Then the polar set $(B')^0 = \{x'' \in (X_s')'\}; \sup\limits_{x' \in B'} |\langle x', x'' \rangle| \leqq 1\}$ of $B'$ is a neighbourhood of $0$ of $(X_s')_s'$, and it is a convex, balanced and absorbing set closed in $\cdot(X_s')'$. Thus $(B')^0 \cap X = {}^0(B')$ is a convex, balanced and absorbing set of $X$. As a (left) polar set, ${}^0(B')$ is closed in the weak topology of $X$, and hence ${}^0(B')$ is closed in the original topology of $X$. Thus ${}^0(B') = (B')^0 \cap X$ is a barrel of $X$, and so it is a neighbourhood of $0$ of $X$. Therefore, the mapping $x \to Jx$ of $X$ into $(X_s')'$ is continuous, because the topology of $(X_s')'$ is defined by a fundamental system of neighbourhoods of $0$ of the form $(B')^0$, where $B'$ ranges over bounded sets of $X'$. 

Let, conversely, $U$ be a convex, balanced and closed neighbourhood of $0$ of $X$. Then, by the preceding section, $U = {}^0 (U^0)$. Thus $JU = JX \cap (U^0)^0$. On the other hand, $U^0$ is a bounded set of $X_s'$, since, for any bounded set $B$ of $X$, there exists an $\alpha > 0$ such that $\alpha B \subseteq U$ and so $(\alpha B)^0 \supseteq U^0$. Hence $(U^0)^0$ is a neighbourhood of $0$ of $(X_s')_s'$. Thus the image $JU$ of the neighbourhood $U$ of $0$ of $X$ is a neighbourhood of $0$ of $JX$ provided with the relative topology of $JX$ as a subset of $(X_s')_s'$. 
\end{proof}


\yossection{3}{Semi-reflexivity and Reflexivity}
\label{sec:AV.3}

\begin{definition}\label{def:AV.3.Definition.1} $A$ locally convex linear topological space $X$ is called semi-reflexive if every continuous linear functional on $X_s'$ is given by 
\begin{equation}\label{eq:AV.3.1}
\langle x, x ^{\prime} \rangle , \text{ with a certain } x \in X.
\end{equation}
\end{definition}

Thus $X$ is semi-reflexive iff 
\begin{equation}\label{eq:AV.3.2}
X _{w} = (X _{s} ^{\prime}) _ {x ^{*}} ^{\prime}.
\end{equation}

\begin{definition}\label{def:AV.3.Definition.2} $A$ locally convex linear topological space $X$ is called reflexive if 
\begin{equation}\label{eq:AV.3.3}
X = (X _{s} ^{\prime}) _{s} ^{\prime}.
\end{equation}
\end{definition}

By Theorem~\ref{thm:AV.3.Theorem.2} in the preceding section, we have 

\begin{proposition}\label{prop:AV.3.Proposition.1} $A$ semi-reflexive space $X$ is reflexive if $X$ is a barrel space. It is also clear, from Definition~\ref{def:AV.3.Definition.2}, that we have 
\end{proposition}

\begin{proposition}\label{prop:AV.3.Proposition.2} The strong dual of a reflexive space is reflexive. 
\end{proposition}

\begin{theorem}\label{thm:AV.3.Theorem.1} $A$ locally convex linear topological space $X$ is semi-reflexive iff every closed, convex, balanced and bounded set of $X$ is compact in the weak topology of $X$. 
\end{theorem}

\begin{proof}
 Let $X$ be semi-reflexive, and $T$ a closed, convex, balanced and bounded set of $X$. Then, by Theorem~\ref{thm:AV.3.Theorem.2} in Section~\hyperref[sec:AV.1]{$1$} of this Appendix, $T = {}^0 (T^0)$. $T$ being a bounded set of $X$, $T^0$ is a neighbourhood of $0$ of $X_s'$. Thus, by Theorem~\ref{thm:AV.3.Theorem.1} in Section~\hyperref[sec:AV.1]{$1$} of this Appendix, $(T^0)^0$ is compact in the weak* topology of $(X_s')'$. Hence, by the semi-reflexivity of $X$, $T = {}^0 (T^0)$ is compact in the weak topology of $X$. 

We next prove the sufficiency part of Theorem~\ref{thm:AV.3.Theorem.1}. Take any $x'' \in (X_s')'$. The strong continuity of $x''$ on $X_s'$ implies that there exists a bounded set $B$ of $X$ such that 
\[
\left| \langle x ^{\prime}, x ^{\prime \prime} \rangle \right| \leqq 1 \text{ whenever } x ^{\prime} \in B ^{0}, \text{ that is, } x ^{\prime \prime} \in (B ^{0}) ^{0}.
\]

We may assume that $B$ is a convex, balanced and closed set of $X$. Thus, by the hypothesis of Theorem~\ref{thm:AV.3.Theorem.1}, $B$ is a compact set in the weak topology of $X$. Hence $B = B^{wa}$ where $B^{wa}$ denotes the closure of $B$ in the weak topology of $X$. Since $X_w$ is embedded in $(X_s')_{w^*}$ as a linear topological subspace, we must have $(B^0)^0 \supseteq B^{wa} = B$. Therefore we have to show that $x''$ is an accumulation point of $B$ in $(X_s')_{w^*}$. Consider the mapping $x \to \varphi(x) = \{\langle x, x_1', \ldots, \langle x, x_n'\rangle\}$ of $X$ into $l^2(n)$, where $x_1', \ldots, x_n' \in X'$. The image $\varphi(B)$ is convex, balanced and compact, since $B$ is convex, balanced and weakly compact. If $\{\langle x_1', x''\rangle, \ldots, \langle x_n', x''\rangle\}$ does not belong to $\varphi(B)$, then, by Mazur's theorem, there would exist a point $\{c_1, \ldots, c_n\} \in l^2(n)$ such that $\sup\limits_{b \in B} \left| \sum\limits_i c_i \langle b, x_i' \rangle \right| \leqq 1$ and $\left( \sum\limits_i c_i \langle x_i', x''\rangle \right) > 1$, proving that $\sum\limits_i c_i x_i' \in B^0$ and $x''$ cannot belong to $B^{00}$. 
\end{proof}

\begin{theorem}\label{thm:AV.3.Theorem.2} $A$ locally convex linear topological space $X$ is reflexive iff it is a barrel space and every closed, convex, balanced and bounded set of $X$ is compact in the weak topology of $X$. In particular, $\mathfrak{D}(\mathbb{R}^{n})$ and $\mathfrak{E}(\mathbb{R}^{n})$ are reflexive. 
\end{theorem}

\begin{proof}
 The sufficiency part is proved already. We shall prove that the first condition of Theorem~\ref{thm:AV.3.Theorem.2} is necessary. 

Let $T$ be a barrel of $X$. We shall prove that $T$ absorbs any bounded set $B$ of $X$ so that $B^0 \supseteq \alpha T^0, \alpha > 0$. $B^0$ being a neighbourhood of $0$ of $X_s'$, we see that $T^0$ is a bounded set of $X_s'$. By the Proposition and Theorem~\ref{thm:AV.3.Theorem.2} in Section~\hyperref[sec:AV.1]{$1$} of this Appendix, we have $T = {}^0(T^0)$. By the hypothesis that $X$ is reflexive, we have ${}^0(T^0) = (T^0)^0$ and so $T = (T^0)^0$. Therefore we have proved that the barrel $T$ is a neighbourhood of $0$ of $X = (X_s')_s'$. Hence $X$ is a barrel space. 

By hypothesis, the closed, convex and balanced set $K = \Conv\left(\bigcup\limits_{|\alpha| \leqq 1} \alpha B\right)^a$ 
is compact in the weak topology of $X$. Here we denote by $\Conv(N)^a$ the closure in $X$ of the convex closure (see p. $28$) $\Conv(N)$ of $N$. Set $Y = \bigcup\limits_{n=1}^{\infty} nK$ and let $p(x)$ be the Minkowski functional of $K$. Then since $K$ is $w$-compact in $Y$, $p(x)$ defines a norm of $Y$. That is, the system $\{\alpha K\}$ with $\alpha > 0$ defines a fundamental system of neighbourhoods of the normed linear space $Y$, and $Y$ is a $B$-space since $K$ is $w$-compact. Hence $Y$ is a barrel space. On the other hand, since $K$ is a bounded set of $X$, the topology of $Y$ defined by the norm $p(x)$ is stronger than the relative topology of $Y$ as a subset of $X$. As a barrel of $X$, $T$ is closed in $X$. Hence $T \cap Y$ is closed in $Y$ with respect to the topology defined by the norm $p(x)$. Therefore $T \cap Y$ is a barrel of the $B$-space $Y$, and so $T \cap Y$ is a neighbourhood of $0$ of the $B$-space $Y$. We have thus proved that $T \cap Y$ and, a fortiori, $T$ both absorb $K \supseteq B$. 
\end{proof}


\yossection{4}{The Eberlein-Shmulyan Theorem}
\label{sec:AV.4}

This theorem is very important in view of its applications. 

\begin{theoremn}[\textsc{Eberlein}-\textsc{Shmulyan}] A $B$-space $X$ is reflexive iff it is locally sequentially weakly compact; that is, $X$ is reflexive iff every strongly bounded sequence of $X$ contains a subsequence which converges weakly to an element of $X$.

For the proof we need two Lemmas: 

\end{theoremn}
\begin{lemma}\label{lem:AV.4.Lemma.1} If the strong dual $X_s'$ of a $B$-space $X$ is separable, then $X$ itself is separable. 
\end{lemma}

\begin{lemma}[S. \textsc{Banach}]\label{lem:AV.4.Lemma.2} $A$ linear subspace $M'$ of the dual space $X'$ of a $B$-space $X$ is weakly* closed iff $M'$ is boundedly weakly* closed; that is, $M'$ is weakly* closed iff $M'$ contains all weak* accumulation points of every strongly bounded subset of $M'$. 
\end{lemma}

\yoslabel{Lemma}{$1$}\label{lem:AV.4.Lemma.1a}is already proved as a Lemma in Chapter~\ref{chap:V}, \ref{sec:V.2}. For Lemma~\ref{lem:AV.4.Lemma.2}, we have only to prove its ``if'' part. It reads as follows. 

\begin{proof}
 (E. \textsc{Hille}-R. S. \textsc{Phillips} \cite{hille97}). We remark that $M'$ is strongly closed by the hypothesis. Let $x_0' \in M'$. Then we can prove that, for each constant $C$ satisfying the condition $0 < C < \inf\limits_{x' \in M'} \|x' - x_0'\|$, there exists an $x_0 \in X$ with $\|x_0\| \leqq 1 / C$ such that 
\begin{equation}\label{eq:AV.4.1}
\langle x _{0}, x _{0} ^{\prime} \rangle = 1 \quad \text{ and } \quad \langle x _{0}, x ^{\prime} \rangle = 0 \quad \text{ for all } \quad x ^{\prime} \in M ^{\prime}. \end{equation}

Thus the strongly closed set $M'$ must contain all of its weak* accumulation points. 


To prove the existence of $x_0$, we choose an increasing sequence of numbers $\{C_n\}$ such that $C_1 = C$ and $\lim\limits_{n \to \infty} C_n = \infty$. Then there exists a finite subset $\sigma_1$ of the unit sphere $S = \{x \in X; \|x\| \leqq 1\}$ such that $\|x' - x_0'\| \leqq C_2$ and $\sup\limits_{x \in \sigma_1} |\langle x, x' \rangle - \langle x, x_0' \rangle| \leqq C_1$ implies $x' \in M'$. 

If not, there would exists, corresponding to each finite subset $\sigma$ of $S$, an $x_{\sigma}' \in M$ such that 
\[
\left\| x _{\sigma} ^{\prime} - x _{0} ^{\prime} \right\| \leqq C _{2} \quad \text{ and } \quad \sup\limits _{x \in \sigma} | \langle x, x _{\sigma} ^{\prime} \rangle - \langle x, x _{0} ^{\prime} \rangle | \leqq C _{1}.
\]

We order the sets $\sigma$ by inclusion relation and denote the weak* closure of the set $\{x_{\sigma}^{\prime};\sigma^{\prime}\geqq \sigma \}$ by $N_{\sigma}^{\prime}$. It is clear that $N_{\sigma}^{\prime}$ enjoys the finite intersection property. On the other hand, since $M^{\prime}$ is boundedly weakly* closed, the Corollary of Theorem~\hyperref[thm:AV.1.Theorem.1]{$1$} in this Appendix, $1$, implies that the set 
\[
M _ {c ^{\prime}} ^{\prime} = \{x ^{\prime} \in M ^{\prime}; \| x ^{\prime} \| \leqq C ^{\prime} \}
\]
is weakly* compact. Hence $N_{\sigma}' \subseteq M_{c'}'$ for $C' = C_2 + \|x_0'|$, and so there exists an $x_1' \in \bigcap\limits_{\sigma} N_{\sigma}' \subseteq M'$. Hence we have $\sup\limits_{x \in S} |\langle x, x_1' \rangle - \langle x, x_0' \rangle| \leqq C_1$ and so $\|x_1' - x_0'| \leqq C_1$, contrary to the hypothesis that $0 < C_1 < \inf\limits_{x' \in M'} \|x' - x_0'|$. 

By a similar argument, we successively prove the existence of a sequence of finite subsets $\sigma_{1},\sigma_{2},\ldots$ of $S$ such that 
\[
\begin{yoslines}
&\| x ^{\prime} - x _{0} ^{\prime} \| \leqq C _{k} \quad \text{ and }  \sup\limits _ {x \in \sigma_{i}} | \langle x, x ^{\prime} \rangle - \langle x, x _{0} ^{\prime} \rangle | \leqq C _{i} (i = 1, 2, \dots , k - 1) \\
&\text{ do not imply } x ^{\prime} \in M ^{\prime}.
\end{yoslines}
\]

Thus, since $\lim\limits_{i\to \infty}C_i = \infty$, we see that $x^{\prime}\bar{\in} M^{\prime}$ if 
\[
\left| \langle x, x ^{\prime} \rangle - \langle x, x _{0} ^{\prime} \rangle \right| \leqq C \quad \text{ for all } \quad x \in (C / C _{j}) \sigma_{j} \quad (j = 1, 2, \dots).
\]

Let $\{x_{n}\}$ be a sequence which successively exhausts the sets $(C / C_j)\sigma_j$ ( $j = 1,2,\ldots$). Then $\lim\limits_{n\to \infty}x_n = 0$ and so $L(x') = \{\langle x_n,x'\rangle \}$ is a bounded linear transformation of $X_s'$ into the $B$-space $(c_0)$. We know that the point $\{\langle x_n,x_0'\rangle \} \in (c_0)$ lies at a distance $>C$ from the linear subspace $L(M')$. Thus, by Corollary of Theorem~\hyperref[thm:IV.6.Theorem.3]{$3$} in Chapter~\ref{chap:IV}, \ref{sec:IV.6}, there exists a continuous linear functional $\{\alpha_{n}\} \in (c_0)' = (l^1)$ such that 
\begin{equation}\left\| \left\{\alpha_{n} \right\} \right\| = \sum\limits_{n=1} ^{\infty} \left| \alpha_{n} \right| \leqq 1 / C, \sum\limits_{n=1} ^{\infty} \alpha_{n} \langle x _{n}, x _{0} ^{\prime} \rangle = 1 \quad \text{ and } \sum\limits_{n=1} ^{\infty} \alpha_{n} \langle x _{n}, x ^{\prime} \rangle = 0 \quad \text{ for all } \quad x ^{\prime} \in M ^{\prime}. \notag\end{equation}

The element $x_0 = \sum\limits_{n=1}^{\infty} \alpha_n x_n$ clearly satisfies condition \eqref{eq:AV.4.1}.
\end{proof}

\begin{corollaryu} Let $\langle x', x_0'' \rangle = F(x')$ be a linear functional defined on the dual space $X'$ of a $B$-space $X$. If $N(F) = N(x_0'') = \{x' \in X'; F(x') = 0\}$ is weakly* closed, then there exists an element $x_0$ such that 
\begin{equation}\label{eq:AV.4.2}
F (x ^{\prime}) = \langle x ^{\prime}, x _{0} ^{\prime \prime} \rangle = \langle x _{0}, x ^{\prime} \rangle \quad \text{ for all } \quad x ^{\prime} \in X ^{\prime}.
\end{equation}
\end{corollaryu}

\begin{proof}
 We may assume that $N(F) \neq X'$. Otherwise we can take $x_{0} = 0$. Let $x_{0}' \in X'$ be such that $F(x_{0}') = 1$. By \eqref{eq:AV.4.1} of the preceding Lemma~\ref{lem:AV.4.Lemma.2}, there exists an $x_{0} \in X$ such that 
\begin{equation}\label{eq:AV.4.3}
\langle x _{0}, x _{0} ^{\prime} \rangle = 1 \quad \text{ and } \quad \langle x _{0}, x ^{\prime} \rangle = 0 \quad \text{ for all } \quad x ^{\prime} \in N (F). \end{equation}

Hence, for any $x' \in X'$, the functional 
\[
x ^{\prime} - F \left(x ^{\prime}\right) x _{0} ^{\prime} = y ^{\prime} \in X ^{\prime}
\]
satisfies $F(y') = 0$, i.e. $y' \in N(F)$. Therefore, by \eqref{eq:AV.4.3}, we obtain \eqref{eq:AV.4.2}. 
\end{proof}

\begin{proof}[Proof of the Theorem.]
``Only if'' part. Let $\{x_{n}\}$ be a sequence of $X$ such that $\|x_{n}\|=1$. The strong closure $X_{0}$ of the subspace spanned by $\{x_{n}\}$ is a separable $B$-space. Being a $B$-space, $X_{0}$ is a barrel space. We shall show that $X_{0}$ is reflexive. Any strongly closed, bounded set $B_{0}$ of $X_{0}$ is also a strongly closed, bounded set of $X$ and hence $B_{0}$ is compact in the weak topology of $X$ by the reflexivity of $X$. But, as a strongly closed linear subspace of $X$, $X_{0}$ is closed in the weak topology of $X$ (see Theorem~\ref{thm:IV.6.Theorem.3} in Chapter~\ref{chap:IV}, \ref{sec:IV.6}). Hence $B_{0}$ is compact in the weak topology of $X_{0}$. Thus $X_{0}$ is reflexive by Theorem~\hyperref[thm:AV.3.Theorem.2]{$2$} in the preceding section. We have thus $X_{0}=((X_{0})_{s}^{\prime})_{s}^{\prime}$. By Lemma~\ref{lem:AV.4.Lemma.1a} above, $(X_{0})_{s}^{\prime}$ is thus separable. Let $\{x_{n}^{\prime}\}$ be strongly dense in $(X_{0})_{s}^{\prime}$. Then the weak topology of $X_{0}$ is defined by an enumerable sequence of semi-norms $p_{m}(x)=|\langle x,x_{m}^{\prime}\rangle|(m=1,2,\ldots)$. Hence it is easy to see that the sequence $\{x_{n}\}$, which is compact in the weak topology of $X_{0}$, is sequentially weakly compact in $X_{0}$ and in $X$ as well. We have only to choose a subsequence $\{x_{n^{\prime}}\}$ of $\{x_{n}\}$ such that finite $\lim\limits_{n\to\infty}\langle x_{n^{\prime}},x_{m}^{\prime}\rangle$ exists for $m$=$1$,$2$, $\ldots$.

``If'' part. Let $M$ be a bounded set of $X$, and assume that every infinite sequence of $M$ contains a subsequence which converges weakly to an element of $X$. We have to show that the closure $\overline{M}$ of $M$ in $X$ in the weak topology of $X$ is weakly compact in $X$. For, then the barrel space $X$ is reflexive by Theorem~\hyperref[thm:AV.3.Theorem.2]{$2$} in the preceding section. 

Since $X_w \subseteq (X_s')'_{w^*}$, we have $\overline{M} = \overline{\overline{M}} \cap X_w$, where $\overline{\overline{M}}$ is the closure of $\overline{M}$ in the weak* topology of $(X_s')'$. Let $S_r'$ be the sphere of $X_s'$ of radius $r > 0$ and centre $0$. By the correspondence 
\begin{equation}\overline {{\overline {{M}}}} \ni m \leftrightarrow \{\langle x ^{\prime}, m \rangle ; \| x ^{\prime} \| \leqq 1 \} \in \prod\limits_ {x ^{\prime} \in S _{1} ^{\prime}} I _ {x ^{\prime}}, \text{ where } I _ {x ^{\prime}} = \{z; | z | \leqq \sup\limits _ {m \in \overline {{M}}} | \langle x ^{\prime}, m \rangle | \}, \notag\end{equation}
$\overline{M}$ may be identified with a closed subset of the topological product $\prod\limits_{x' \in S_i'} I_{x'}$. By Tychonov's theorem, $\prod\limits_{x' \in S_i'} I_{x'}$ is compact and hence $\overline{\overline{M}}$ is compact in the weak* topology of $(X_s')'$. Hence we have only to show $\overline{\overline{M}} \subseteq X_w$. 

Let $x_0'' \in (X_s)'$ be an accumulation point of the set $\overline{M}$ in the weak* topology of $(X_s)'$. To prove that $x_0'' \in X_w$ we have only to show that the set $N(x_0'') = \{x' \in X'; \langle x', x_0'' \rangle = 0\}$ is weakly* closed. For, then, by the above Corollary, there exists an $x_0 \in X$ such that $\langle x', x_0'' \rangle = \langle x_0, x' \rangle$ for all $x' \in X'$. We shall first show that 
for every finite set $x_1', x_2', \ldots, x_n'$ of $X'$, there exists a $z \in \overline{M}$ such that $\langle x_j', x_0'' \rangle = \langle z, x_j' \rangle (j = 1, 2, \ldots, n)$. 

The proof is as follows. Since $x_0''$ is in the weak* closure of $\overline{M}$, there is an element $z_m \in \overline{M}$ such that 
\begin{equation}\label{eq:AV.4.4}
\left| \left\langle z _{m}, x _{j} ^{\prime} \right\rangle - \left\langle x _{j}, x _{0} ^{\prime \prime} \right\rangle \right| \leqq 1 / m (j = 1, 2, \dots , n).
\end{equation}

By hypothesis there exists a subsequence of $\{z_m\}$ which converges weakly to an element $z \in X$ and so $z \in \overline{M}$, since the sequential weak closure of $\overline{M}$ is contained in $\overline{M}$. We have thus \eqref{eq:AV.4.4}. 

Now, by Lemma~\ref{lem:AV.4.Lemma.2}, $N(x_0'')$ is weakly* closed if, for every $r > 0$, the set $N(x_0'') \cap S_r'$ is weakly* closed. Let $y_0'$ be in the weak* closure of $N(x_0'') \cap S_1'$. We have to show that $y_0' \in N(x_0'') \cap S_1'$. To this purpose, we choose an arbitrary $\varepsilon > 0$ and construct three sequences $\{z_n\} \subseteq \overline{M}$, $\{x_n\} \subseteq M$ and $\{y_n'\} \subseteq N(x_0'') \cap S_1'$ as follows: By \eqref{eq:AV.4.4} we can choose a, $z_1 \in \overline{M}$ such that $\langle z_1, y_0' \rangle = \langle y_0', x_0'' \rangle$. $z_1$ being in the weak closure of $M$, there exists an $x_1 \in M$ such that $|\langle x_1, y_0' \rangle - \langle z_1, y_0' \rangle| \leqq \varepsilon/4$. $y_0'$ being in the weak* closure of $N(x_0'') \cap S_1'$, there exists a $y_1' \in N(x_0'') \cap S_1'$ such that $|\langle x_1, y_1' \rangle - \langle x_1, y_0' \rangle| \leqq \varepsilon/4$. Repeating the argument and remembering \eqref{eq:AV.4.4}, we obtain $\{z_n\} \subseteq \overline{M}$, $\{x_n\} \subseteq M$ and $\{y_n'\} \subseteq N(x_0'') \cap S_1'$ such that 
\begin{equation}\label{eq:AV.4.5}
\begin{yoslines}
&\langle z _{1}, y _{0} ^{\prime} \rangle = \langle y _{0} ^{\prime}, x _{0} ^{\prime \prime} \rangle , \\
&\langle z _{n}, y _{m} ^{\prime} \rangle = \langle y _{m} ^{\prime}, x _{0} ^{\prime \prime} \rangle = 0  (m = 1, 2, \dots , n - 1), \\
&| \langle x _{n}, y _{m} ^{\prime} \rangle - \langle z _{n}, y _{m} ^{\prime} \rangle | \leqq \varepsilon / 4  (m = 0, 1, \dots , n - 1), | \langle x _{i}, y _{n} ^{\prime} \rangle - \langle x _{i}, y _{0} ^{\prime} \rangle | \leqq \varepsilon / 4  (i = 1, 2, \dots , n).
\end{yoslines}
\end{equation}

Thus we have 
\begin{equation}\label{eq:AV.4.6}
\left| \langle y _{0} ^{\prime}, x _{0} ^{\prime \prime} \rangle - \langle x _{i}, y _{n} ^{\prime} \rangle \right| \leqq \varepsilon / 4 + \varepsilon / 4 = \varepsilon / 2 (i = 1, 2, \dots , n).
\end{equation}

Since $\{x_{n}\} \subseteq M$, there exists a subsequence of $\{x_{n}\}$ which converges weakly to an element $x \in \overline{M}$. Without losing the generality, we may assume that the sequence $\{x_{n}\}$ itself converges weakly to $x \in \overline{M}$. Hence, from \eqref{eq:AV.4.5}, $|\langle x, y_{m}' \rangle| \leqq \varepsilon / 4$. From $w - \lim\limits_{n \to \infty} x_{n} = x$ and Mazur's theorem $2$ in Chapter~\ref{chap:V}, \ref{sec:V.1}, there exists a convex combination $u = \sum\limits_{j=1}^{n} \alpha_{j} x_{j} \left( \alpha_{j} \geqq 0, \sum\limits_{j=1}^{n} \alpha_{j} = 1 \right)$ such that $\|x - u\| \leqq \varepsilon / 4$. Therefore, by \eqref{eq:AV.4.6}, 
\[
\left| \langle y _{0} ^{\prime}, x _{0} ^{\prime \prime} \rangle - \langle u, y _{n} ^{\prime} \rangle \right| \leqq \sum\limits_{j=1} ^{n} \alpha_{j} \left| \langle y _{0} ^{\prime}, x _{0} ^{\prime \prime} \rangle - \langle x _{j}, y _{n} ^{\prime} \rangle \right| \leqq \varepsilon / 2,
\]
and hence 
\[
\begin{aligned}
&\left| \langle y _{0} ^{\prime}, x _{0} ^{\prime \prime} \rangle \right|  \leqq \left| \langle y _{0} ^{\prime}, x _{0} ^{\prime \prime} \rangle - \langle u, y _{n} ^{\prime} \rangle \right| + \left| \langle u, y _{n} ^{\prime} \rangle - \langle x, y _{n} ^{\prime} \rangle \right| + \left| \langle x, y _{n} ^{\prime} \rangle \right| \leqq \varepsilon / 2 + \| u - x \| \| y _{n} ^{\prime} \| + \varepsilon / 4 \leqq \varepsilon .
\end{aligned}
\]

As $\varepsilon$ was arbitrary, we see that $\langle y_0', x_0'' \rangle = 0$ and so $y_0' \in N(x_0'')$. Combined with the fact that $S_1'$ is weakly* closed, we finally obtain $y_0' \in N(x_0'') \cap S_1'$. 
\end{proof}

\begin{remarku} As for the weak topologies and duality in $B$-spaces, there is an extensive literature. See, e.g., the references in N. \textsc{Dunford}-J. \textsc{Schwartz} \cite{dunford44}. Sections $1$, $2$ and $3$ of this Appendix are adapted and modified from N. \textsc{Bourbaki} \cite{bourbaki26} and A. \textsc{Grothendieck} \cite{grothendieck87}. It is remarkable that necessary tools for proving this far reaching theorem of \textsc{Eberlein} \cite{dunford52}-\textsc{Shmulyan} \cite{shmulyan302} are found, in one form or other, in the book of S. \textsc{Banach} \cite{banach9}. 
\end{remarku}
